% Generated draft: requires mathematical and source review.
\documentclass[11pt]{article}
\usepackage[margin=25mm]{geometry}
\usepackage{fontspec}
\setmainfont{Latin Modern Roman}
\usepackage{amsmath,amssymb,mathtools,unicode-math}
\setmathfont{Latin Modern Math}
\usepackage{enumitem,needspace,xurl,hyperref}
\setlength{\parindent}{0pt}
\setlength{\parskip}{5pt}
\setlength{\emergencystretch}{3em}
\newcommand{\N}{\mathbb{N}}
\newcommand{\Z}{\mathbb{Z}}
\newcommand{\Q}{\mathbb{Q}}
\newcommand{\R}{\mathbb{R}}
\newcommand{\C}{\mathbb{C}}
\newcommand{\F}{\mathbb{F}}
\newcommand{\A}{\mathbb{A}}
\newcommand{\PP}{\mathbb{P}}
\newcommand{\E}{\mathbb{E}}
\newcommand{\eps}{\varepsilon}
\newcommand{\dd}{\,\mathrm{d}}
\begin{document}
\section*{UnsolvedMath ID: 30006316}
\textbf{Source problem code:} catalogue entry
\subsection{Definitions and mathematical statement}
Let \(F\) be a nonarchimedean local field, let \(\mathcal O\) be its ring of integers, and let \(\mathbf F_q\) be its residue field, of characteristic \(p\). Let \(G\) be a split connected reductive group over \(F\), and put \(K=G(\mathcal O)\), subject to specifying an integral model. Let \(X\) be a spherical \(G\)-variety: geometrically, a Borel subgroup has a dense orbit. Assume that \(T^*X\) is hyperspherical in the sense intended by the cited relative Langlands duality framework. The supplied record does not give the definition of this additional condition or all geometric hypotheses needed for it.

Choose the integral models and the \(G(F)\)-eigenmeasure \(\nu_X\) intended by the source. An eigenmeasure means that translation scales the measure by a character; the character and the corresponding normalization of the action are not supplied. The damaged normalization display appears to prescribe
\[
 \nu_X\bigl(X(\mathcal O)\bigr)
 =\frac{\lvert X(\mathbf F_q)\rvert}{\lvert G(\mathbf F_q)\rvert}
 q^{\dim G-\dim X}.
\]
This parsing requires confirmation. In particular, the reductions and integral points in this formula require compatible integral models. Set
\[
 \mathcal S_X=C_c^\infty(X(F))^K,
 \qquad b_X=\mathbf 1_{X(\mathcal O)},
 \qquad
 \langle f,h\rangle_X=\int_{X(F)}f(x)\overline{h(x)}\,\dd\nu_X(x).
\]
Here \(C_c^\infty\) denotes compactly supported locally constant complex-valued functions. The intended hypotheses must ensure that \(b_X\in\mathcal S_X\). Let \(\mathcal H(G(F),K)\) be the convolution algebra of compactly supported \(K\)-bi-invariant functions. Its Haar normalization and its action on \(\mathcal S_X\) must be chosen according to the source's requirement that the action be unitary; the record does not specify the corresponding adjoint convention.

For a conditional formulation of the proposed spectral side, write \(H=G_X^\vee\) for the intended dual group, \(A_X\subset H\) for its indicated torus, \(W_X\) for the corresponding Weyl group, and \(H_c\subset H(\C)\) for the indicated compact form. The source also specifies the compact torus as the intersection of the indicated dual torus with this compact form. Some superscripts in these data are displaced in the transcription, so possible twists or modified groups remain unresolved. With that qualification, the intended quotient is
\[
 Z=H\mathbin{//}H\simeq A_X\mathbin{//}W_X,
 \qquad
 \C[H\mathbin{//}H]=\C[H]^H,
\]
where \(H\) acts on itself by conjugation. Let \(V_X\) denote the graded self-dual representation of the intended dual group referred to in source reference \([1,\S4.5]\). Its construction is not recoverable from the damaged expression involving \(S_X\), Lie algebras, an intersection, and an orthogonal-complement symbol in the supplied record.

Let \(\pi:H_c\to Z(\C)\) be the quotient map. The proposed measure is the pushforward
\[
 \mu=\pi_*\left(
 \frac{\dd g}{\det(1-q^{-1/2}g\mid V_X)}
 \right),
\]
where \(\dd g\) is Haar measure on \(H_c\). Its scale, the meaning of the determinant for the graded representation, and the hypotheses ensuring a positive measure and finite polynomial norms require confirmation. For \(a,c\in\C[Z]\), the intended pairing is
\[
 \langle a,c\rangle_\mu
 =\int_{Z(\C)}a(z)\overline{c(z)}\,\dd\mu(z).
\]

\textbf{Numerical unramified Plancherel conjecture, conditional on the unresolved source data.}
For every \((F,G,X)\) satisfying the intended hypotheses and normalizations, there is a complex-linear isomorphism
\[
 \Phi:\mathcal S_X\xrightarrow{\sim}\C[Z]
\]
such that, for all \(f,h\in\mathcal S_X\) and \(T\in\mathcal H(G(F),K)\),
\[
 \Phi(b_X)=1,\qquad
 \langle\Phi(f),\Phi(h)\rangle_\mu=\langle f,h\rangle_X,
 \qquad
 \Phi(Tf)=\rho_X(T)\Phi(f).
\]
Here \(\rho_X\) is the intended spectral homomorphism from the spherical Hecke algebra to \(\C[Z]\); its construction is not specified in the record. The target consists of regular functions on the affine quotient, rather than an unspecified space of all functions. This is a statement about the displayed function spaces with pairings; the source does not explicitly replace them by Hilbert completions.

The original transcription also interleaves example-specific Hecke formulas and a unit-circle integral yielding \(q^{-n/2}=\langle T_n b_X,b_X\rangle\). The example, strata, indices, and normalization are not adequately specified, so these computations cannot be imposed as general hypotheses. Finally, the source writes the coefficient field as \(k=\Q_\ell\simeq\C\), with \(\ell\ne p\). This cannot be a literal field isomorphism; whether an algebraic-closure symbol or another convention was lost is unresolved.
\subsection{Short English statement}
The proposed conjecture identifies compactly supported unramified functions on a spherical variety with regular functions on a dual quotient, preserving the basic vector and inner product and turning Hecke operators into multiplication. Its predicted measure is expressed using a dual representation. Several hypotheses, normalizations, and dual-data conventions remain unresolved in the supplied transcription, so this is a conditional statement-only draft requiring mathematical review.
\subsection{Sources}
\begin{itemize}
\item Local UnsolvedMath registry, permanent ID \(30006316\), exact source problem code \texttt{catalogue entry}, title \emph{The Numerical Unramified Plancherel Conjecture}. The supplied record belongs to the Oberwolfach Reports Open Problems collection and gives proposed year 2025; no individual proposer is supplied.
\item \emph{Arbeitsgemeinschaft: Relative Langlands Duality} (2025), Oberwolfach Reports, pp.~811--882, DOI \texttt{10.4171/owr/2025/18}. Supplied report and source link: \url{https://doi.org/10.4171/owr/2025/18}. The extracted text labels the statement ``Conjecture 2.1'' and attaches the reference \([1,\text{Proposition }9.2.1]\); the supplied record does not explicitly identify that numbered bibliography entry. It also refers to \([1,\S4.5]\) and, within interleaved example material, \([1,(9.3)]\).
\item Sakellaridis, \emph{Spherical functions on spherical varieties}: \url{https://arxiv.org/abs/0905.4244}.
\item Ben-Zvi, Sakellaridis, and Venkatesh, \emph{Relative Langlands Duality}: \url{https://www.math.ias.edu/~akshay/research/BZSVpaperV1.pdf}.
\item Supplied literature assessment, checked exactly \texttt{2026-08-22}: the record labels the problem partially solved and states that the numerical unramified Plancherel statement remains conjectural in the general hyperspherical setting, that the report contains matching computations in several examples, and that Sakellaridis proves classical spherical cases. This is supplied dated curation, not independently verified here; the scope of the cited results and time-sensitive status require independent verification.
\item The registry labels the statement \texttt{corrected\_verified} and says that the official report was checked, the question isolated from surrounding exposition, and notation repaired. These are supplied provenance claims, not a source visit or mathematical verification performed for this draft.
\end{itemize}
\end{document}
